\documentclass[conference]{IEEEtran}

\usepackage{graphicx}
\usepackage{amsmath,amssymb}
\usepackage{float}
\usepackage{booktabs}
\usepackage{multirow}

\begin{document}

\title{
    \textbf{Benchmark Report: 523.xalancbmk\_r} \\
    \text{Register Spill Analysis on RISC-V via gem5}
}

\author{
    \IEEEauthorblockN{Lütfullah Burak Kaya}
    \IEEEauthorblockA{
        Ozyegin University \\
        burak.kaya.50711@ozu.edu.tr
    }
    \and
    \IEEEauthorblockN{Ismail Akturk}
    \IEEEauthorblockA{
        Ozyegin University \\
        ismail.akturk@ozyegin.edu.tr
    }
}

\newcommand{\LongDate}{%
    \number\day\space
    \ifcase\month
    \or January\or February\or March\or April\or May\or June%
    \or July\or August\or September\or October\or November\or December%
    \fi
    \space \number\year
}

\twocolumn[
\begin{@twocolumnfalse}
\maketitle
\begin{center}{\small \LongDate}\end{center}
\vspace{1em}
\end{@twocolumnfalse}
]

% ================================================

\begin{abstract}
This report presents the register spill characterization of the
\texttt{523.xalancbmk\_r} benchmark from SPEC CPU2017, executed on a
RISC-V (RV64GC) target using the gem5 TimingSimpleCPU simulator.
A custom SpillDetector was integrated into gem5 to count store-then-load
patterns on the stack within a user-defined Region of Interest (ROI).
The test dataset (\texttt{test.xml} transformed via \texttt{xalanc.xsl})
was used as the input. Results show a spill rate of \textbf{4.64\%}
within the ROI, corresponding to 18.4 million detected spill events
across 396 million ROI instructions. XPath expression evaluation and
DOM tree traversal maintain a large set of simultaneously live node
pointers and context objects, creating persistent register pressure
throughout the XSLT transformation pipeline.
\end{abstract}

% =========================================================
\section{Benchmark Overview}

\texttt{523.xalancbmk\_r} is the SPEC CPU2017 rate variant of
Xalan-C++ 1.10, an XSLT processor that applies XSL Transformation
stylesheets to XML input documents. The benchmark exercises XML
parsing (via Xerces-C++), XSL stylesheet compilation, XPath expression
evaluation, and DOM tree traversal with namespace resolution. It is
one of the most object-oriented workloads in the SPEC suite, with a
large C++ class hierarchy and heavy use of virtual dispatch throughout
the transformation pipeline.

The RISC-V binary (\texttt{xalancbmk\_r.riscv}) was compiled with gem5
ROI markers inserted around the \texttt{xsltMain()} call in
\texttt{src/XalanExe.cpp}:

\begin{itemize}
    \item \texttt{m5\_work\_begin(0,0)} --- before \texttt{xsltMain()}
    \item \texttt{m5\_work\_end(0,0)}   --- after  \texttt{xsltMain()}
\end{itemize}

\texttt{xsltMain()} encompasses the entire transformation pipeline:
input parsing, stylesheet compilation, XPath evaluation, and output
serialisation.

\subsection{Input Datasets}

SPEC CPU2017 provides three dataset sizes for \texttt{xalancbmk\_r},
differing in the size and complexity of the XML input document
(Table~\ref{tab:inputs}).

\begin{table}[H]
\centering
\caption{xalancbmk\_r Input Dataset Sizes}
\label{tab:inputs}
\begin{tabular}{lll}
\toprule
\textbf{Dataset} & \textbf{XML input} & \textbf{Description} \\
\midrule
test    & test.xml (28\,KB)   & Small SPEC benchmark table \\
train   & Various XML files   & Medium-sized documents \\
refrate & t5.xml ($\approx$100\,MB) & Large XHTML document, repeated \\
\bottomrule
\end{tabular}
\end{table}

The \textbf{test} dataset was selected for this study. The input
\texttt{test.xml} is a 28\,KB XML document containing SPEC benchmark
reference data, transformed by the 16\,KB \texttt{xalanc.xsl}
stylesheet into an XHTML table. The transformation exercises XPath
predicates, template matching, and recursive element processing
representative of the full Xalan-C++ code paths.

% =========================================================
\section{Simulation Setup}

\subsection{gem5 Configuration}

\begin{itemize}
    \item \textbf{CPU model:}   TimingSimpleCPU
    \item \textbf{ISA:}         RISC-V RV64GC
    \item \textbf{Caches:}      enabled (default L1 I/D + L2)
    \item \textbf{Memory:}      4\,GB simulated physical RAM
    \item \textbf{Spill mode:}  summary only (\texttt{spill\_verbose=False})
\end{itemize}

The full gem5 invocation was:

\begin{verbatim}
./build/RISCV/gem5.opt
  --outdir=benchmarks/xalancbmk/m5out_test
  configs/deprecated/example/se.py
  --cmd=benchmarks/xalancbmk/xalancbmk_r.riscv
  --options="-v benchmarks/xalancbmk/test.xml
             benchmarks/xalancbmk/xalanc.xsl"
  --cpu-type=TimingSimpleCPU
  --caches
  --mem-size=4GB
\end{verbatim}

% =========================================================
\section{Simulation Results}

The simulation completed successfully. The output was a correctly
formed XHTML document containing the SPEC benchmark reference table,
confirming the transformation completed without errors.

\subsection{Pre-ROI Context}

Before \texttt{xsltMain()} was entered, the binary executed
Xalan-C++ initialisation: loading the Xerces XML parser, initialising
the XSLT processor, and registering built-in XPath functions.
Table~\ref{tab:preroi} shows the global counters at ROI entry.

\begin{table}[H]
\centering
\caption{Global Counters at ROI Entry}
\label{tab:preroi}
\begin{tabular}{lr}
\toprule
\textbf{Metric} & \textbf{Value} \\
\midrule
Total instructions (pre-ROI) & 2,689,452 \\
Total spills detected (pre-ROI) & 90,538 \\
Pre-ROI spill rate & 3.37\% \\
\bottomrule
\end{tabular}
\end{table}

\subsection{ROI Spill Statistics}

Table~\ref{tab:roi} presents the spill statistics collected within
the ROI.

\begin{table}[H]
\centering
\caption{ROI Spill Statistics --- 523.xalancbmk\_r (test.xml)}
\label{tab:roi}
\begin{tabular}{lr}
\toprule
\textbf{Metric} & \textbf{Value} \\
\midrule
ROI instructions         & 395,919,950 \\
ROI stores               &  38,990,477 \\
ROI loads                & 106,032,294 \\
\midrule
\textbf{ROI spills}      & \textbf{18,381,250} \\
\textbf{Spill rate}      & \textbf{4.6427\%} \\
\bottomrule
\end{tabular}
\end{table}

\subsection{Timing}

\begin{table}[H]
\centering
\caption{Simulation Timing}
\label{tab:timing}
\begin{tabular}{lr}
\toprule
\textbf{Metric} & \textbf{Value} \\
\midrule
Simulated time           & 0.968\,s \\
Host wall time           & 280.72\,s ($\approx$4.7\,min) \\
Total simulated instructions & 401,232,382 \\
\bottomrule
\end{tabular}
\end{table}

% =========================================================
\section{Analysis}

\subsection{Spill Rate Interpretation}

The measured spill rate of \textbf{4.64\%} places \texttt{xalancbmk\_r}
fifth among all evaluated SPEC CPU2017 benchmarks
(Table~\ref{tab:compare}), between \texttt{526.blender\_r} (6.22\%)
and \texttt{541.leela\_r} (4.44\%).

\begin{table}[H]
\centering
\caption{Cross-Benchmark Spill Rate Comparison}
\label{tab:compare}
\begin{tabular}{lrr}
\toprule
\textbf{Benchmark} & \textbf{ROI Spills} & \textbf{Spill Rate} \\
\midrule
507.cactuBSSN\_r & 3,782,099,438 & 7.8756\% \\
500.perlbench\_r &     1,070,055 & 7.2838\% \\
502.gcc\_r       &     1,108,838 & 6.9185\% \\
526.blender\_r   &    61,845,869 & 6.2193\% \\
523.xalancbmk\_r &    18,381,250 & \textbf{4.6427\%} \\
541.leela\_r     & 1,158,236,271 & 4.4438\% \\
519.lbm\_r       &   274,130,440 & 4.2130\% \\
531.deepsjeng\_r & 1,928,903,656 & 4.1749\% \\
544.nab\_r       &   151,310,616 & 2.2612\% \\
538.imagick\_r   &     2,424,899 & 2.0601\% \\
505.mcf\_r       &   446,113,136 & 1.8263\% \\
508.namd\_r      &   327,718,955 & 1.0269\% \\
\bottomrule
\end{tabular}
\end{table}

\subsection{Store vs.\ Load Balance}

The load count (106.0M) is $2.72\times$ higher than the store count
(39.0M). This strongly load-dominated ratio characterises XPath
evaluation: each template match requires traversing the DOM tree by
following node pointers, loading the node type, name, and attribute
list, and evaluating XPath predicates against the node content.
The same DOM nodes are frequently revisited during recursive template
instantiation, generating many loads for each store (the stores
correspond to writing output nodes into the result tree).

\subsection{Spill Fraction of Loads}

Within the ROI, the ratio of spills to total loads is:
\[
\frac{18{,}381{,}250}{106{,}032{,}294} \approx 17.3\%
\]
Approximately 1 in 6 load instructions is a spill reload.
The Xalan-C++ transformation pipeline passes a rich context object
through each recursive template invocation: the current source node,
the current result tree fragment, the XPath evaluation context, the
variable binding frame, and the namespace resolution context. With
RISC-V providing 32 general-purpose registers and the heavy use of
virtual dispatch adding call-frame overhead, these context pointers
are continuously spilled to and reloaded from the stack.

\subsection{Pre-ROI Spill Density}

The pre-ROI spill rate of 3.37\% ($90{,}538$ spills over
$2{,}689{,}452$ instructions) reflects the Xalan-C++ and Xerces
initialisation phase: loading the XML parser, registering XSLT
extension functions, and initialising the XPath function library.
This rate is comparable to \texttt{gcc\_r}'s pre-ROI rate (3.34\%),
suggesting that complex C++ framework initialisation consistently
produces moderate register pressure.

\subsection{ROI Coverage}

The ROI accounts for $395{,}919{,}950$ out of $401{,}232{,}382$
total simulated instructions:
\[
\frac{395{,}919{,}950}{401{,}232{,}382} \approx 98.7\%
\]
The near-perfect ROI coverage confirms that \texttt{xsltMain()}
encompasses virtually all benchmark computation, with only 1.3\%
of instructions outside the ROI for program startup and teardown.

\subsection{XSLT-Driven Register Pressure}

The 4.64\% spill rate is structurally driven by Xalan-C++'s
object-oriented XSLT engine. During each template instantiation step,
the following objects must be kept simultaneously live:
\begin{itemize}
    \item The current DOM source node (\texttt{XalanNode*})
    \item The XPath evaluation context (\texttt{XPathExecutionContext})
    \item The current result tree builder (\texttt{FormatterListener*})
    \item The active variable binding stack (\texttt{XalanQName})
    \item The namespace resolver (\texttt{PrefixResolver*})
\end{itemize}
Each virtual function call in the dispatch chain saves and restores
these pointers, as the RISC-V backend cannot keep all of them in
registers simultaneously. The 2.72:1 load-to-store ratio directly
reflects the read-heavy nature of DOM traversal, where nodes are
read many times per write to the output tree.

% =========================================================
\section{Conclusion}

The \texttt{523.xalancbmk\_r} benchmark exhibits a register spill rate
of 4.64\% within the \texttt{xsltMain()} ROI on RISC-V RV64GC, the
fifth highest among all evaluated SPEC CPU2017 benchmarks. With 18.4
million spill events across 396 million ROI instructions and a 98.7\%
ROI coverage, the benchmark is well-characterised by this run. The
2.72:1 load-to-store ratio reflects the read-heavy DOM traversal
pattern of XPath evaluation, and the 17.3\% spill-to-load fraction
arises from the rich set of live context pointers that the Xalan-C++
template instantiation engine must maintain simultaneously. These
results place \texttt{523.xalancbmk\_r} firmly in the moderate-to-high
register-pressure tier of the evaluated SPEC suite.

\end{document}
